\documentclass[12pt,a4paper]{article}
\usepackage[utf8]{vntex, inputenc}
\usepackage{amsmath, amssymb, anyfontsize, array, color,  enumerate, enumitem, fancyhdr, float, graphicx, hyperref, lastpage, listings, mathtools, mips, natbib, pifont, tikz}
\usepackage[open,openlevel=0]{bookmark}
\usepackage[framemethod=tikz]{mdframed}

\usetikzlibrary{calc}
\newcommand\HRule{\rule{12cm}{1pt}}

\pagestyle{fancy}
\fancyhead{} % clear all header fields
\fancyhead[L]{
  \color{blue}
  \begin{tabular}{rl}
    \begin{picture}(25,15)(0,0)
    \put(0,-8){\includegraphics[width=8mm, height=8mm]{BK.png}}
    \end{picture}
    \begin{tabular}{l}
      \textbf{\bf \ttfamily Ho Chi Minh University of Technology}\\
      \textbf{\bf \ttfamily Faculty of Computer Science \& Engineering}
    \end{tabular}
  \end{tabular}
}
\fancyhead[R]{
  \begin{tabular}{l}
    \tiny \bf \\
    \tiny \bf
  \end{tabular}
}
\fancyfoot{} % clear all footer fields
\fancyfoot[L]{\scriptsize \ttfamily Computer Architecture Week 5}
\fancyfoot[R]{\scriptsize \ttfamily Page {\thepage}/\pageref{LastPage}}
\renewcommand{\headrulewidth}{0.3pt}
\renewcommand{\footrulewidth}{0.3pt}

\definecolor{dkgreen}{rgb}{0,0.6,0}
\definecolor{gray}{rgb}{0.5,0.5,0.5}
\definecolor{mauve}{rgb}{0.58,0,0.82}

\lstset{frame=tb,
  language=[MIPS]Assembler,
  aboveskip=3mm,
  belowskip=3mm,
  showstringspaces=false,
  columns=flexible,
  basicstyle={\small\ttfamily},
  numbers=none,
  numberstyle=\tiny\color{gray},
  keywordstyle=\color{blue},
  commentstyle=\color{dkgreen},
  stringstyle=\color{mauve},
  breaklines=true,
  breakatwhitespace=true,
  tabsize=3,
  numbers=left,
  stepnumber=1,
  numbersep=1pt,
  firstnumber=1,
  numberfirstline=true
}

\begin{document}
\begin{titlepage}
  \begin{tikzpicture}[remember picture, overlay]
    \draw[line width = 2pt] ($(current page.north west) + (1in,-1in)$) rectangle ($(current page.south east) + (-1in,1in)$);
  \end{tikzpicture}

  \begin{center}

    % Upper part of the page
    \textbf{\fontsize{12pt}{1pt}\selectfont HO CHI MINH CITY UNIVERSITY OF TECHNOLOGY}\\[0.5cm]
    {\fontsize{13pt}{1pt}\selectfont Faculty of Computer Science \& Engineering}\\[0.5cm]
    \begin{figure}[H]
      \centering
      \includegraphics[width=1.7in,height=1.7in]{BK.png}
    \end{figure}

    % Title
    \HRule\\[0.5cm]
    { \textbf{\fontsize{25pt}{1pt}\selectfont Computer Architecture Week 6}}\\[0.4cm]

    \HRule\\[0.8cm]
    \begin{minipage}{0.545\textwidth}
      \begin{flushleft}
        \textbf{Authors:}\\
        \begin{tabular}{l l}
          Lưu Nguyễn Hoàng Minh & 1952845 \\
          Vũ Anh Nhi            & 1952380 \\
          Nguyễn Chính Khôi     & 1952793 \\
          Nguyễn Hoàng          & 1952255 \\
        \end{tabular}
      \end{flushleft}
    \end{minipage}
    \begin{minipage}{0.4\textwidth}
      \begin{flushright}
        \textbf{Class:}\\
        CO2008\_CC02\\
        \textbf{Lecturer:}\\
        Kiều Đỗ Nguyên Bình\\

      \end{flushright}
    \end{minipage}

    \vfill

    % Bottom of the page
    \vspace{2cm}
    {\large} %{\large \today}
  \end{center}
\end{titlepage}


\pdfbookmark[0]{Question 1}{ques1}
\section*{Question 1:}
Given the following program with MIPS assembly.
\begin{mdframed}[hidealllines=true,backgroundcolor=magenta!10]
  \begin{lstlisting}
    .data
      s_a: .float 1.1
      s_b: .float 2.2

    .text
      la      $t0, s_a
      lwc1    $f0, 0($t0)
      la      $t0, s_b
      lwc1    $f1, 0($t0)
      add.s   $f2, $f1, $f0

      li      $v0, 10
      syscall
  \end{lstlisting}
\end{mdframed}
Students run the program and answer the following question
\begin{enumerate}
  \item Explain the function of lwc1, add.s instructions. \\
        - lwc1: Load word into coprocessor 1 \\
        - add.s: Floating point addition single precision
  \item Can we use lw to load value for \$f0 register? Explain. \\
        - No, because floating point registers are separate from ordinary registers
  \item Is the result that is stored in \$f2 3.3? If it is not, explain. \\
        - No. In base-2, the only prime factor is 2, so we can only cleanly express fractions whose denominator has only 2 as a prime factor.
        So 0.1 and 0.2 (11/10 and 6/5) are repeating decimals in base-2. \\
        When we perform math on repeating decimals, we end up with leftovers which carry over when converting base-2 to base-10.
  \item Analyze the values that are stored in \$f0, \$f1, \$f2 based on IEEE 754 standard. \\
        - \$f0: 0\(\times\)3f8ccccd = 0011 1111 1000 1100 1100 1100 1100 1101
        \begin{align*}
          Sign           = & \;0                                                         \\
          Exponent       = & \;127                                                       \\
          Fraction       = & \;2^{-4} + 2^{-5} + 2^{-8} + 2^{-9} + 2^{-12}  +            \\
                           & \;2^{-13} + 2^{-16} + 2^{-17} + 2^{-20} + 2^{-21} + 2^{-23} \\
          Decimal\;value = & \;{(-1)}^{0} \times 2^{(127 - 127)} \times (1 + Fraction)   \\
          =                & \;1.1000000238
        \end{align*}
        - \$f1: 0\(\times\)400ccccd = 0100 0000 0000 1100 1100 1100 1100 1101
        \begin{align*}
          Sign           = & \;0                                                         \\
          Exponent       = & \;128                                                       \\
          Fraction       = & \;2^{-4} + 2^{-5} + 2^{-8} + 2^{-9} + 2^{-12}  +            \\
                           & \;2^{-13} + 2^{-16} + 2^{-17} + 2^{-20} + 2^{-21} + 2^{-23} \\
          Decimal\;value = & \;{(-1)}^{0} \times 2^{(128 - 127)} \times (1 + Fraction)   \\
          =                & \;2.2000000477
        \end{align*}
        - \$f2: 0\(\times\)40533334 = 0100 0000 0101 0011 0011 0011 0011 0100
        \begin{align*}
          Sign           = & \;0                                                          \\
          Exponent       = & \;128                                                        \\
          Fraction       = & \;2^{-1} + 2^{-3} + 2^{-6} + 2^{-7} + 2^{-10}  +             \\
                           & \;2^{-11} + 2^{-14} + 2^{-15}  + 2^{-18} + 2^{-19} + 2^{-21} \\
          Decimal\;value = & \;{(-1)}^{0} \times 2^{(128 - 127)} \times (1 + Fraction)    \\
          =                & \;3.300000191
        \end{align*}
\end{enumerate}


\pdfbookmark[0]{Question 2}{ques2}
\section*{Question 2:}
Given the following program with MIPS assembly.
\begin{mdframed}[hidealllines=true,backgroundcolor=magenta!10]
  \begin{lstlisting}
    .data
      d_a: .double 1.1
      d_b: .double 2.2

    .text
      la      $t0, d_a
      ldc1    $f0, 0($t0)
      la      $t0, d_b
      ldc1    $f1, 0($t0)
      add.d   $f2, $f1, $f0

      li      $v0, 10
      syscall
  \end{lstlisting}
\end{mdframed}
Students run the program and answer the following question
\begin{enumerate}
  \item List the differences between this program and the program in Question 1. \\
        - This program stores data with double precision floating point registers. \\
        - This program doesn't run because the first register must be even numbered.
  \item Is the result that is stored in \$f2 3.3? If it is not, explain.
        - No. In base-2, the only prime factor is 2, so we can only cleanly express fractions whose denominator has only 2 as a prime factor.
        So 0.1 and 0.2 (11/10 and 6/5) are repeating decimals in base-2. \\
        When we perform math on repeating decimals, we end up with leftovers which carry over when converting base-2 to base-10.
  \item Analyze the values that are stored in \$f0, \$f1, \$f2 based on IEEE 754 standard.
  \item Correct the program to add 2 double number 1.5 and 2.25.
        \begin{mdframed}[hidealllines=true,backgroundcolor=magenta!10]
          \begin{lstlisting}
    .data
      d_a: .double 1.5
      d_b: .double 2.25

    .text
      la      $t0, d_a
      ldc1    $f0, 0($t0)
      la      $t0, d_b
      ldc1    $f2, 0($t0)
      add.d   $f4, $f2, $f0

      li      $v0, 10
      syscall
          \end{lstlisting}
        \end{mdframed}
\end{enumerate}

\pdfbookmark[0]{Question 3}{ques3}
\section*{Question 3:}
Implement a procedure that convert input degree in Fahrenheit to Celsius.
If the degree in Celsius higher than 99.5, print ``Warning'' to terminal, otherwise print ``Safe'' to terminal.

\begin{mdframed}[hidealllines=true,backgroundcolor=magenta!10]
  \begin{lstlisting}
    # TODO
    # Answer is attached in .zip file
  \end{lstlisting}
\end{mdframed}

\end{document}
